\section*{Capítulo \ref{ch:g6:pure-substances-and-mixtures} --- Sustancias puras y mezclas}

\begin{solution}{exo:g6:pure-substances-and-mixtures:1}
\omterm{def:g6:pure-substances-and-mixtures:pure-substance}{Sustancias puras}: el agua destilada, el bloque de hierro puro, el
cristal de azúcar. \omterm{def:g6:pure-substances-and-mixtures:mixture}{Mezclas}: el \omterm{def:g4:air-a-mixture-of-gases:air}{aire}, el agua del mar, la leche.
\end{solution}

\begin{solution}{exo:g6:pure-substances-and-mixtures:2}
Homogéneas: el agua salada, el agua del grifo. Heterogéneas: el granito,
el aceite sobre agua, el muesli.
\end{solution}

\begin{solution}{exo:g6:pure-substances-and-mixtures:3}
Una clase concreta de sustancia, con su propio nombre y sus propias
propiedades: por ejemplo el agua, la sal, el etanol (también el azúcar,
el oxígeno, el hierro).
\end{solution}

\begin{solution}{exo:g6:pure-substances-and-mixtures:4}
Una temperatura de ebullición fija sugiere una \omterm{def:g6:pure-substances-and-mixtures:pure-substance}{sustancia pura};
\qty{78}{\celsius} es la temperatura de ebullición del etanol.
\end{solution}

\begin{solution}{exo:g6:pure-substances-and-mixtures:5}
Contiene muchas \omterm{def:g6:pure-substances-and-mixtures:species}{especies químicas} (agua, azúcares, ácidos, vitaminas,
aromas): es una \omterm{def:g6:pure-substances-and-mixtures:mixture}{mezcla}. ``Puro'' en el brik solo quiere decir que no se
le ha añadido nada.
\end{solution}

\begin{solution}{exo:g6:pure-substances-and-mixtures:6}
No: una \omterm{def:g6:pure-substances-and-mixtures:pure-substance}{sustancia pura} hierve a una temperatura fija. Es una \omterm{def:g6:pure-substances-and-mixtures:mixture}{mezcla},
probablemente agua con algo disuelto, como agua salada.
\end{solution}

\begin{solution}{exo:g6:pure-substances-and-mixtures:7}
Un filtro podría separar el agua turbia (los granos se quedan en el
papel) y el zumo con pulpa (la pulpa se queda). No puede separar el
agua salada (la sal está disuelta) ni el aceite del agua (los dos
líquidos atraviesan el filtro; se separan vertiendo la capa de arriba).
\end{solution}

\begin{solution}{exo:g6:pure-substances-and-mixtures:8}
$78.7 \div 10 = \qty{7.87}{g}$ por \qty{1}{cm^3}: hierro.
\end{solution}

\begin{solution}{exo:g6:pure-substances-and-mixtures:9}
$\frac{357}{478} \approx 0.747$: alrededor del \qty{75}{\percent}.
\end{solution}

\begin{solution}{exo:g6:pure-substances-and-mixtures:10}
$2 \times 35 = \qty{70}{g}$ de sales. Porcentaje: $\frac{35}{1000} =
\qty{3.5}{\percent}$.
\end{solution}

\begin{solution}{exo:g6:pure-substances-and-mixtures:11}
Medir sus temperaturas de ebullición (cada una debe mantenerse fija
mientras hierve, y las dos deben ser iguales) y pesar el mismo volumen
de cada uno (masas iguales). Dos constantes iguales, ambas fijas, apuntan
a la misma \omterm{def:g6:pure-substances-and-mixtures:pure-substance}{sustancia pura}.
\end{solution}

\begin{solution}{exo:g6:pure-substances-and-mixtures:12}
$\frac{10}{90 + 10} = \qty{10}{\percent}$ de etanol. Es una \omterm{def:g6:pure-substances-and-mixtures:mixture}{mezcla}: no
tiene una temperatura de ebullición fija y no aparece en la tabla, que
solo enumera \omterm{def:g6:pure-substances-and-mixtures:pure-substance}{sustancias puras}.
\end{solution}

\begin{solution}{pb:g6:pure-substances-and-mixtures:1}
\textbf{1.} No. Una \omterm{def:g6:pure-substances-and-mixtures:pure-substance}{sustancia pura} y una \omterm{def:g6:pure-substances-and-mixtures:homogeneous}{mezcla homogénea} pueden
parecer exactamente iguales: hay que medir algo.

\textbf{2.} Homogénea: es transparente y sus partes no se ven.

\textbf{3.} Una constante: una temperatura de ebullición que se mantiene
fija mientras el líquido hierve (o la masa de un volumen conocido,
comparada con una tabla).

\textbf{4.} A y B, cuyas temperaturas de ebullición se mantienen fijas,
son \omterm{def:g6:pure-substances-and-mixtures:pure-substance}{sustancias puras}. C es una \omterm{def:g6:pure-substances-and-mixtures:mixture}{mezcla}.

\textbf{5.} A hierve a \qty{100}{\celsius}: agua. B hierve a
\qty{78}{\celsius}: etanol.

\textbf{6.} A: $49.8 \div 50.0 = \qty{0.996}{g}$. B: $39.5 \div 50.0 =
\qty{0.790}{g}$. C: $51.2 \div 50.0 = \qty{1.024}{g}$.

\textbf{7.} Sí: unos \qty{1.0}{g} para el agua y \qty{0.79}{g} para el
etanol.

\textbf{8.} A medida que el agua se va al hervir, la sal se queda, así
que el líquido que queda contiene cada vez más sal; entonces hierve a
una temperatura más alta.

\textbf{9.} El líquido C contiene un sólido disuelto: es una disolución,
una \omterm{def:g6:pure-substances-and-mixtures:mixture}{mezcla}, como ya mostraba su ebullición.

\textbf{10.} $100 \times 1.024 = \qty{102.4}{g}$; $\frac{3.5}{102.4}
\approx 0.034$: alrededor del \qty{3.4}{\percent} de la masa es sal.

\textbf{11.} $\frac{35}{1000} = \qty{3.5}{\percent}$: muy cerca del
resultado de la pregunta 10.

\textbf{12.} \qty{3.5}{g} en \qty{100}{mL}, así que $3.5 \times 10 =
\qty{35}{g}$ de sal por litro. Por su contenido en sal y por cómo
hierve, C es el agua del mar; A es el agua destilada y B, el etanol.
\end{solution}
